% Original CUT and box trip proofs; no modern switching criterion is used.
\section{Original CUT ports, contractions and cyclic boxes}

\begin{annexitem}{annex:p07.girard-cut.erasure}
 {thm:p07-girard-cut-erasure}{All marker positions and the exact orbit correspondence}
\begin{proof}
At position L the marker changes the direct crossing \(X^-\mapsto Y^+\)
to \(X^-\mapsto\kappa^-\mapsto\kappa^+\mapsto Y^+\);
the other crossing is unchanged. At R it subdivides the other crossing
\(Y^-\mapsto X^+\) by those same two gates.
Thus erasure gives precisely the direct successor at both positions.
Every orbit containing a marker also contains an ordinary premise gate;
there is no marker-only orbit. Deleting the two inserted states from
each orbit therefore gives a bijection of marked and direct orbits.
Insertion reverses it. Repeat independently at every cut.
For a fixed ordinary setting all marker choices correspond to the
same direct orbits, proving the stated equivalence of full correctness.
\end{proof}
\end{annexitem}

\begin{annexitem}{annex:p07.girard-cut.sequentialization}
 {thm:p07-girard-cut-sequentialization}{Tagged roots and every proof-construction case}
\begin{proof}
Replace each marker by a distinct tagged occurrence of the actual
formula \(X\otimes X^\perp\), produced by its two cut premises,
and keep that occurrence terminal. The successor pattern is exactly
the marked tensor pattern, including both positions and terminal return.
The resulting ordinary cut-free structure is correct.
Original cut-free sequentialization gives its proof, with each tagged
root produced by its designated tensor inference.

That root has no consumer in the graph, so it is never principal in
a later inference. Replace its designated tensor inference by cut,
and erase only that tagged occurrence from every later side context.
The two principal premise occurrences and their disjoint proofs remain.
All other inferences keep their actual principal occurrences.
Equal formula labels elsewhere are not erased.
Do this at each tag to obtain the desired sequent proof with cuts,
reconstructing exactly the original cut incidences.

Conversely, an axiom, exchange, par or tensor proof has the original
long trips by the already proved construction theorem.
At cut, its two disjoint premise nets each have one orbit and a specified
terminal occurrence. Replace their two terminal returns by the two
direct cross-connections. This joins the two orbits into one.
The erasure lemma then inserts its marker in either position.
Induction proves correctness for every ordinary and marker setting.
No semantic soundness or cut-preservation theorem is used here.
\end{proof}
\end{annexitem}

\begin{annexitem}{annex:p07.girard-cut.principal-preservation}
 {thm:p07-girard-cut-principal-preservation}{Four exhaustive slices and both multiplicative orientations}
\begin{proof}
Fix an arbitrary setting of every ordinary link retained by the reduct.
Carry it to the original, and first take
\(U=B\otimes C,\ V=B'\bindnasrepma C'\), setting both removed
producer positions to L. These positions are free preimages, not
positions of the reduct, so this choice restricts no reduct setting.
By erasure the old direct permutation has one orbit.

Reinsert the selected CUT marker temporarily. Its two positions both
have one orbit. The three-successor calculation for an original tensor
therefore applies to its port pattern: cutting the three changing arcs
gives disjoint paths \(U^+\ldots U^-\), \(V^+\ldots V^-\),
and the marker return, which has only its two marker gates.
Thus the first two paths partition all ordinary gates.
This use of the three-path calculation depends only on the actual
successors and both full-orbit settings, not on a formula label for CUT.

The original tensor calculation at U, with its root gates as endpoints,
decomposes its premise interval into
\(P_C:C^+\ldots C^-\) followed by \(P_B:B^+\ldots B^-\).
The par calculation at V decomposes its interval into
\(Q_1:B'^+\ldots C'^-\) followed by \(Q_2:C'^+\ldots B'^-\).
These four disjoint paths cover every ordinary gate except \(U^\pm,V^\pm\).
The old direct orbit, starting at \(B^-\), is exactly
\[
 B^-,U^-,V^+,Q_1,Q_2,V^-,U^+,P_C,P_B,B^-,
\]
where a path symbol includes its displayed endpoints.
This is just the source's local successor clauses joined to the
two passage-order normal forms.

After contraction the direct cuts give
\[
 B^-,Q_1,P_C,Q_2,P_B,B^-.
\]
Indeed their crossings are \(B^-\mapsto B'^+\),
\(B'^-\mapsto B^+\), \(C^-\mapsto C'^+\),
\(C'^-\mapsto C^+\). The four unchanged paths are visited once,
so the reduct has one full direct orbit.
Insert each new marker at either position; the erasure lemma proves
correctness at all these positions. All retained positions were arbitrary.

For the other orientation, the first producer is par and its opposite
is tensor. Exchange the two symmetric cut operands and rename the
paired premises \(B,C\) as \(B',C'\), and conversely.
No connective's left/right ports are exchanged. The same four-path
calculation gives the cuts on the same corresponding pairs.
This covers both principal orientations and every reduct setting.
\end{proof}
\end{annexitem}

\begin{annexitem}{annex:p07.girard-cut.axiom-preservation}
 {thm:p07-girard-cut-axiom-preservation}{Both axiom sides, both choices and the prohibited loop}
\begin{proof}
Let the selected axiom output be X, its companion Z, and the opposite
cut premise Y. First suppose \(Z,Y\) are distinct.
In the direct permutation the exact local chains are
\[
 Z^+\mapsto X^-\mapsto Y^+,\qquad
 Y^-\mapsto X^+\mapsto Z^-.
\]
The glued occurrence W has plus gate inherited from \(Y^+\),
whose producer is retained, and minus gate from \(Z^-\),
whose consumer or terminal boundary is retained.
Erase \(Z^+,X^-,Y^-,X^+\) by shortening those two chains;
retarget their incoming incidences to the retained endpoints.
The resulting successor is precisely the glued net's direct successor.
At least \(Y^+,Z^-\) remain. Shortening this one cyclic orbit leaves
one orbit over every remaining gate. Retained switch choices were
arbitrary, and arbitrary remaining marker positions follow by erasure.

If \(Z=Y\), the cut ties together that same axiom's two outputs.
Its direct orbits are \(X^+,Y^-,X^+\) and \(Y^+,X^-,Y^+\).
Marker insertion subdivides only one and cannot join them.
Such an object is not correct, excluding exactly the size/gluing trap
in the original source.

Selecting an axiom on the other side exchanges the two roles and
gives the same chain-shortening argument, not a new premise.
If both sides have distinct axiom producers, either application leaves
one axiom joining the two companion occurrences, with the same labels
and retained outside incidences. The results are isomorphic.
Thus both choices and all mirrored cases preserve correctness.
\end{proof}
\end{annexitem}

\begin{annexitem}{annex:p07.girard-cut.normalization}
 {thm:p07-girard-cut-normalization}{The source measure and every irreducible-cut possibility}
\begin{proof}
A principal step removes two ordinary roots and one marker and inserts
two markers. Thus \(p'=p-2,\ c'=c+1\), so \(s'=s-1\).
An axiom step replaces three ordinary occurrences X,Z,Y by one glued
occurrence and removes one marker. Hence \(p'=p-2,\ c'=c-1\),
so \(s'=s-3\). The same-axiom case that could invalidate this
count was excluded by explicit short trips in the preservation proof.

Each intermediate net remains correct by the two preservation theorems.
Every formula occurrence in a closed PN0 net has an axiom, tensor or
par producer. If either cut premise has an axiom producer, its axiom
step applies. Otherwise dual outer constructors force a tensor opposite
a par, so the principal step applies. Thus no irreducible cut remains.
The positive integer s strictly decreases, forbidding an infinite path.
The reduct remains nonempty: a principal step keeps all four premise
occurrences, and an axiom step keeps the glued occurrence.
Hence its final size is at least one. At least one size unit is lost
per step, giving at most \(s-1\) steps. This is the original measure,
with no box or open-hypothesis case added to it.
\end{proof}
\end{annexitem}

\begin{annexitem}{annex:p07.girard-cut.confluence}
 {thm:p07-girard-cut-confluence}{All incidences that can overlap}
\begin{proof}
Different cuts cannot share a tensor/par producer, whose unique output
has at most one consumer. If their internal redexes are disjoint,
the surgeries commute: each changes only its specified incidences,
retaining the endpoints needed by the other, even when adjacent.

If two different cuts share an axiom, they use its two different outputs.
Delete either cut and that axiom. One cut remains between the two
former opposite premises. The two outcomes have the same labels and
outside endpoints, up to renaming the internal cut and its marker and
the cut's symmetric order. At one cut, both axiom choices similarly
leave one axiom between the two companion outputs, with the same
retained ordinary boundary identities. The prohibited same-axiom
loop cannot occur in a correct net. A nonaxiom principal redex has
only its corresponding-premise rule. These exhaust overlaps, proving
local confluence for the actual source surgeries.

For full confluence induct on the decreasing s.
For reductions starting \(N\to U,N\to V\), local confluence supplies
a common descendant W. The smaller-measure induction inside U joins
its chosen endpoint with a descendant of W; induction inside V does
the same for the other endpoint. Induction below W then joins these
descendants. Zero-step cases are immediate. Thus all finite reductions
join. By normalization, maximal endpoints are cut-free; their common
reduct can only be an isomorphic copy. Ordinary retained boundary
identities and labels are preserved; generated internal markers can
be renamed. No modern correctness or semantic soundness theorem has
entered this overlap argument.
\end{proof}
\end{annexitem}

\begin{annexitem}{annex:p07.girard-box.slices}
 {thm:p07-girard-box-slices}{All cyclic choices force self-returning exterior paths}
\begin{proof}
Choose one cyclic position at the given box and cut its n successor arcs
out of the one full orbit. The remainder is n disjoint directed paths
covering all gates, with heads \(F_i^-\) and tails \(F_j^+\).
They are unchanged by changing only this box's position.
If the path from \(F_i^-\) ends at \(F_j^+\) with \(j\ne i\),
choose an n-cycle \(\rho\) with \(\rho(j)=i\).
Such a cycle exists: list j,i and then all remaining indices in a
cyclic order. Reconnecting closes precisely that path into an orbit.
The other paths are nonempty, so it is a proper short orbit, impossible.
Therefore every path ends at its own door's plus gate.
For \(n=1\) this follows directly from the sole head and tail.
This uses the whole class of allowed n-cycles, not all permutations.
\end{proof}
\end{annexitem}

\begin{annexitem}{annex:p07.girard-box.empire}
 {thm:p07-girard-box-empire}{Door, same-slice, different-slice and singleton cases}
\begin{proof}
Fix every other position and use the proved exterior slices
\(S_i:F_i^-\ldots F_i^+\).
Changing the box position changes only the cyclic order of these
unchanged paths.

If the distinguished tensor premise A is itself door \(F_k\), its
interval from \(A^+\) to \(A^-\) starts at the end of \(S_k\),
traverses every other slice, and ends at the start of \(S_k\).
It contains both gates of every door, including k, in every cyclic
position. This proves the assertion in this boundary case.

Otherwise \(A^+,A^-\) are non-endpoint gates in these slices.
If they are in the same slice and plus is earlier than minus,
their interval is internal and contains no complete door endpoint pair.
If minus is earlier, the interval wraps through all other slices and
the two ends of this slice, and contains every door pair.
If the gates are in different slices \(S_i,S_j\), the start door's
minus endpoint and end door's plus endpoint are outside the interval.
For any third door k, choose a cyclic order with i immediately
followed by j. Then \(S_k\) is not traversed before the minus endpoint,
so its door pair is also excluded. Such a position is allowed;
for \(n=2\) there is no third door to consider.
Thus no door can have both gates universally in the A-interval in
the different-slice case. For \(n=1\), only the already covered
same-slice possibilities exist.

If one door belongs to the universal empire, the excluding cases
are impossible for every fixed choice of outside positions.
Only a distinguished-door case or a wrapping same-slice case remains,
and each includes all door pairs at every cyclic position.
As the outside choices were arbitrary, all doors belong to the empire.
\end{proof}
\end{annexitem}

\begin{annexitem}{annex:p07.girard-box.sequentialization}
 {thm:p07-girard-box-sequentialization}{The complete original empire route with whole boxes}
\begin{proof}
For the forward direction, a box axiom with all doors terminal alternates
\(F_i^+\mapsto F_{\rho(i)}^-\mapsto F_{\rho(i)}^+\).
An n-cycle visits all \(2n\) gates once. Ordinary dual axioms have the
original four-gate orbit. Par inserts its two root gates into one
terminal return; tensor joins two disjoint orbits and inserts its root
gates; exchange changes none. Induction proves the forward direction
at every box position and every ordinary position.

We supply the converse's empire dependencies rather than assume that
an ordinary axiom can simply be replaced by an arbitrary box.
Fix a distinguished tensor premise A. At any other tensor or par,
fix all outside positions, including the box cycles. Removing the
three changing successor arcs gives exactly the same three-path
endpoint permutations as the original local-block calculation:
tensor forces the identity matching; par forces the transposition
of its two premise names. This calculation needs only the one-cycle
hypothesis and the local clauses, so its six-case proof still applies.
The abstract interval-position proof is unchanged. It gives all four
tensor/par closure facts, including root=A, the distinguished-tensor
exclusion, and universal rather than one-setting membership.
The distinguished premise belongs to its empire; the two premise
empires are disjoint and exclude their tensor root.
An ordinary axiom's companion closure follows from its adjacent gates.
A box's producer closure is exactly the all-door empire lemma proved
above. Thus every producer of an empire occurrence has all its outputs
and premises within that empire.

Prescribe reflecting positions at pars with exactly one inside premise.
The trip-realization exit proof now has one additional case:
at a box door's plus gate it jumps to another door's minus gate,
which remains inside by all-door closure, for every cyclic position.
The other plus, minus, terminal and distinguished-boundary cases are
the original tensor/par cases. Hence the interval from \(A^+\) to
\(A^-\) equals all gates of its empire, with every non-boundary setting,
including all box cycles, free.

Producer ancestry outside a box uses only connective premise edges,
since a box has no exterior premises. Its free positions can choose
a plus-only path from A to either child of an upstream tensor.
That child's minus gate cannot occur earlier on the plus-only initial
path. Its full interval is therefore contained in the realized A
interval, proving inclusion of its universal empire. This also includes
a tensor immediately producing A.

Choose an inclusion-maximal tensor empire union when there are no
terminal pars. If its root is not terminal, follow its consumers to
a terminal connective; formula complexity increases at each tensor/par
step, so finiteness ends the chain. It is a tensor.
The just-proved ancestor inclusion puts the old union in the
appropriate larger premise empire, strictly so because the old tensor
root is in that larger empire but outside its own two empires.
This contradicts maximality, proving terminality.

If that terminal tensor does not cover the structure by its two
empires and root, compatible reflecting prescriptions would realize
both empires simultaneously unless there is a crossing par.
With no such crossing the root's two adjacent terminal gates and
the two premise intervals partition every gate, forcing covering.
Thus noncovering gives a par with one premise in each empire.
Its root is in neither, by par producer closure. Follow its consumers
to a terminal tensor, and realize the appropriate larger premise
empire with a plus-only producer path to that par.
All its producer choices are free inside the larger empire.
First choose the par's left position. Its passage order is
\(C^+,D^-,D^+,C^-\) between its root gates.
Disjointness of the old premise intervals forces the old D-side
interval into the arc \(C^+\ldots C^-\), inside the larger realized
interval. Universal membership then puts that entire old empire in
the larger empire. At the right position, with the same free ancestry
path, the order is \(D^+,C^-,C^+,D^-\); the symmetric interval argument
puts the other old empire there. These are two separately realized
settings. Their union is strictly enlarged by the crossing root.
This again contradicts maximality. No modern splitting theorem is used.

At a covering tensor each retained box is wholly in one restriction:
one door inside puts all its doors inside, and disjointness prevents
sharing the same box across both. Axiom and connective producer
closure keep every producer complete. Consumer closure keeps every
foreign tensor complete; a crossing par would have a root in neither
restriction, contrary to covering. Only the selected tensor is removed.
There are no reflecting boundary pars left. Any positions of a
restriction, including each whole box's original n-cycle, extend to
positions of the full net. The realization exit argument visits
exactly that restriction's gates before the distinguished minus gate;
its new terminal return closes one full orbit. Both restrictions are
therefore nonempty correct boxed structures, without splitting a box.

Now induct on the number of exterior links. A terminal par is removed
by deleting its root's two consecutive gates and exposing its premises;
the reduced trip is one orbit at every setting. Otherwise the proved
covering tensor gives the two smaller correct restrictions.
Apply induction and the corresponding sequent rule.
With no connective, multiple axiom/box producers would have separate
orbits, so there is just one ordinary axiom or one entire box.
It is the matching designated axiom leaf. This proves cut-free
sequentialization, using every original box occurrence once.

For CUT markers, replace them temporarily by tagged terminal tensor
formulas as in the earlier original-cut proof. The boxed cut-free
result applies. Erase each tagged root from later side contexts and
replace its designated tensor inference by cut; its producer and
terminal status prevent it from being a later principal occurrence.
The forward cut construction joins the two disjoint premise orbits
and inserts the dummy ports. This proves the stated extension, not
normalization against arbitrary box leaves.
\end{proof}
\end{annexitem}

\begin{annexitem}{annex:p07.girard-box.plugging}
 {thm:p07-girard-box-plugging}{The induced cyclic permutation and every inner gate}
\begin{proof}
Rename all internal non-door occurrences and markers fresh before gluing.
Fix an arbitrary internal setting. Its correct trip is one cycle.
Remove the n terminal returns \(F_i^-\mapsto F_i^+\) of the
ordinary doors, but not any internal CUT returns.
This leaves n disjoint paths covering all inner gates, from \(F_i^+\)
to \(F_j^-\). Let \(\rho(i)=j\).
Restoring the terminal returns reconnects these paths by the permutation
\(\rho\); since the whole inner trip is one cycle, \(\rho\) is an
n-cycle. Thus it is an allowed exterior box position.

Fix the outside settings and choose this \(\rho\) at that box.
The correct exterior has one orbit. Replace each jump
\(F_i^+\mapsto F_{\rho(i)}^-\) by its corresponding internal path.
The paths have disjoint interiors, contain every inner non-door gate,
and meet the exterior only at their correctly matched endpoints.
Expanding these edges consequently gives one orbit over exactly all
plugged gates. Formula labels and producer/consumer identities match
at every door. Internal and external settings were arbitrary, proving
the assertion. No claim that every n-cycle is realized by the given
interior is needed; exterior checking covers whichever allowed cycle
that interior actually induces.
\end{proof}
\end{annexitem}

\begin{annexitem}{annex:p07.girard-box.unary-insertion}
 {thm:p07-girard-unary-insertion}{The two-gate insertion and inverse}
\begin{proof}
Originally terminal F has the successor \(F^-\mapsto F^+\).
After appending the source unary root G it is replaced by
\(F^-\mapsto G^-\mapsto G^+\mapsto F^+\):
the middle arrow is G's terminal return, and the other arrows are
the actual unary port clauses. Exactly two fresh gates subdivide
that edge, so the original one orbit remains one over every new gate.
There is no additional switch. Removing a terminal unary root
shortens precisely this chain, proving the inverse for every setting.
The typing constructor must be one of the specified source unary
rules; the argument supplies no arbitrary unary axiom or full box
normalization.
\end{proof}
\end{annexitem}

